% ======================================================================
%  udt — plik startowy: SKRYPT (tryb script)
%
%  Dłuższe teksty: skrypty, prace zaliczeniowe, notatki. Jedna kolumna,
%  blok tytułowy, spis treści, numerowane rozdziały, bibliografia.
%  Skopiuj ten plik razem z udt-core.sty i udt-script.sty.
%  Środowiska twierdzeń i matematyka działają tak samo jak w trybie
%  report — szczegółowy opis w streszczenie.tex; rysunki, tabele,
%  jednostki i bibliografia — w sprawozdanie.tex.
%
%  Kompilacja:  pdflatex → biber → pdflatex → pdflatex
%  Wersja rozszerzona (wszystkie warianty) z tego samego źródła:
%      pdflatex -jobname=skrypt-ext skrypt
% ======================================================================

% ----------------------------------------------------------------------
%  OPCJE
%    plain — środowiska twierdzeń bez kolorowych bloków,
%    chem  — chemia: \ce{...} (mhchem) oraz
%              \std            — indeks stanu standardowego (⦵),
%              \Eo             — potencjał standardowy E⦵,
%              \redox{Ox}{Red} — para redoks Ox/Red,
%    full  — włącza wszystkie warianty (jak nazwa zadania z -ext).
% ----------------------------------------------------------------------
\documentclass[a4paper, polish, 11pt]{scrartcl}
\usepackage[chem]{udt-script}
\addbibresource{literatura.bib}

% ----------------------------------------------------------------------
%  WARIANTY — treść warunkowa
%  Jedno źródło, kilka wersji dokumentu. Wariant to środowisko, którego
%  treść składa się tylko wtedy, gdy wariant jest włączony; inaczej
%  znika, jakby go nie było w źródle.
%    \newvariant{nazwa1, nazwa2}  — deklaracja (preambuła),
%    \showvariant{nazwa1}         — włączenie (preambuła, dowolna kolejność),
%    \begin{nazwa1} ... \end{nazwa1} — treść (wcięcia dozwolone).
%  Wariant extended jest predefiniowany. Opcja full albo nazwa zadania
%  kończąca się na -ext włącza wszystkie warianty, więc
%      pdflatex skrypt   i   pdflatex -jobname=skrypt-ext skrypt
%  dają skrypt.pdf i skrypt-ext.pdf z jednego pliku.
%  Włączona treść ma pełną szerokość i cienką kreskę na lewym
%  marginesie; może przechodzić przez strony. Rysunki i tabele w środku
%  stoją tam, gdzie je wpisano ([H]).
%  Ograniczenia ukrytej treści: zrównoważone klamry, bez verbatim.
%  Ukryta treść nie dostaje numerów — numeracja wersji się różni,
%  a \ref spoza wariantu do etykiety w jego wnętrzu daje „??”.
%  Dlatego w wariantach używa się środowisk z gwiazdką (remark*, …).
% ----------------------------------------------------------------------
\newvariant{dodatek}
\showvariant{dodatek}   % zakomentuj, aby ukryć treść wariantu dodatek

% ----------------------------------------------------------------------
%  POLA BLOKU TYTUŁOWEGO
%  Tryb script drukuje: \institution, \course, \group, \place, \date,
%  \title, \subtitle, \author, \studentid, \supervisor.
%  Nie drukuje: \doctype, \exercise, \experimentdate (należą do trybu
%  report). Kilku autorów: \author{A \and B}, \studentid{1 \and 2}.
%  Tytuł i autorzy trafiają też do metadanych PDF.
% ----------------------------------------------------------------------
\institution{Politechnika Gdańska, Wydział Fizyki Technicznej i Matematyki Stosowanej}
\course{Nazwa przedmiotu}
\group{Fizyka Techniczna, Fizyka Stosowana, rok III}
\title{Tytuł skryptu}
\subtitle{Podtytuł}
\author{Imię Nazwisko}
\studentid{000000}
\supervisor{dr Imię Nazwisko}
\place{Gdańsk}
\date{\today}

\begin{document}
\maketitle

% abstract i preface — oba opcjonalne, przed spisem treści.
\begin{abstract}
Krótkie streszczenie: o czym jest tekst i co z niego wynika.
\end{abstract}

\begin{preface}
Przedmowa: dla kogo jest tekst, co trzeba umieć wcześniej, jak go czytać.
\end{preface}

% Spis treści: rozdziały i podrozdziały (dwa poziomy).
\tableofcontents

% ----------------------------------------------------------------------
%  STRUKTURA
%  \part — nagłówek części w tekście (bez nowej strony); numeracja
%          rozdziałów biegnie przez części bez zerowania,
%  \section, \subsection — numerowane i w spisie treści,
%  \subsubsection — numerowane, poza spisem treści.
%  Równania, rysunki, tabele i środowiska twierdzeń numerowane w obrębie
%  rozdziału: 1.1, 1.2, … Pagina: numer i tytuł bieżącego rozdziału.
%  Odsyłacze: \cref{sec:pojecia} → „rozdział 1”, \cref{eq:I} →
%  „równanie (1.1)”, \cref{tw:steiner} → „twierdzenie 1.2”.
% ----------------------------------------------------------------------
\part{Podstawy}

\section{Pojęcia}\label{sec:pojecia}

\subsection{Definicje}\label{sec:definicje}

% Kolejność dedukcyjna: najpierw przesłanki i mechanizm, definicja
% lub twierdzenie jako wniosek.
Rozłożenie masy bryły względem osi obrotu opisuje jedna liczba.

\begin{definition}[Moment bezwładności]\label{def:I}
Momentem bezwładności bryły względem osi o wersorze $\vu{n}$ nazywa się
\begin{equation}
	I = \int h^2 \d m, \qquad h^2 = \rho^2 - (\vu{n} \vdot \vb{\rho})^2
	\label{eq:I}
\end{equation}
gdzie $h$ jest odległością elementu masy $\d m$ od osi.
\end{definition}

\subsection{Twierdzenia}\label{sec:twierdzenia}

\begin{theorem}[Twierdzenie Steinera (o osiach równoległych)]\label{tw:steiner}
Moment bezwładności względem osi odległej o $d$ od równoległej osi przez środek masy wynosi $I = I_0 + md^2$.
\end{theorem}

\begin{proof}
Wystarczy rozpisać $h^2$ względem środka masy; wyrazy liniowe znikają z definicji środka masy.
\end{proof}

\begin{lemma}
Dla bryły płaskiej $I_z = I_x + I_y$.
\end{lemma}

\begin{corollary}
Moment bezwładności jest najmniejszy względem osi przechodzącej przez środek masy (\cref{tw:steiner}).
\end{corollary}

\begin{fact}
Moment bezwładności jest addytywny względem podziału bryły na części.
\end{fact}

\begin{law}[II zasada dynamiki ruchu obrotowego]
$M = I\varepsilon$ dla obrotu wokół ustalonej osi.
\end{law}

\begin{example}
Dla pręta o długości $l$ względem osi przez koniec: $I = \frac{1}{3}ml^2$.
\end{example}

\begin{remark}
Środowiska z gwiazdką (\texttt{remark*} itd.) wyglądają tak samo, ale nie mają numeru.
\end{remark}

\begin{aside}[Kim był Steiner?]
Jakob Steiner był szwajcarskim geometrą XIX wieku.
\end{aside}

\part{Rozszerzenia}

\section{Warianty i chemia}\label{sec:warianty}

% Wariant dodatek — widoczny, bo włączony w preambule \showvariant.
\begin{dodatek}
\begin{remark*}
Ta uwaga należy do wariantu \texttt{dodatek}: znika po zakomentowaniu \texttt{\textbackslash showvariant\{dodatek\}}.
\end{remark*}
\end{dodatek}

% Wariant extended — widoczny tylko z opcją full albo w kompilacji -ext.
\begin{extended}
\begin{aside}[Wariant extended]
Ten blok pojawia się wyłącznie w wersji rozszerzonej.
\end{aside}
\end{extended}

% Chemia (opcja chem). Wzory i reakcje: \ce{...}.
Ogniwo Daniella: \ce{Zn + Cu^2+ -> Zn^2+ + Cu}, siła elektromotoryczna $\Eo = \Eo_{\redox{Cu^2+}{Cu}} - \Eo_{\redox{Zn^2+}{Zn}} = \qty{1.10}{\volt}$, entalpia standardowa $\Delta H\std$.

% Język można zmienić w środku dokumentu; napisy pakietu (Rys., Definicja,
% nazwy w \cref) i format liczb przełączają się razem z nim. Klasa musi
% znać oba języki: \documentclass[english, polish]{scrartcl} (ostatni
% jest domyślny); fragment w innym języku:
%   \begin{otherlanguage}{english} ... \end{otherlanguage}

Literatura uzupełniająca: \cite{taylor,przyklad}.

\printbibliography

\end{document}
